\documentclass[11pt]{article}

% --- packages ---
\usepackage[utf8]{inputenc}
\usepackage[T1]{fontenc}
\usepackage{amsmath, amssymb, amsthm, amsfonts}
\usepackage{mathtools}
\usepackage{booktabs}
\usepackage{enumitem}
\usepackage[margin=1in]{geometry}
\usepackage[hidelinks]{hyperref}
\usepackage{verbatim}
\usepackage{xcolor}
\usepackage{microtype}

% --- theorem environments (shared counter per section) ---
\theoremstyle{definition}
\newtheorem{definition}{Definition}[section]
\newtheorem{condition}[definition]{Condition}
\newtheorem{remark}[definition]{Remark}
\newtheorem{observation}[definition]{Observation}

\theoremstyle{plain}
\newtheorem{theorem}[definition]{Theorem}
\newtheorem{lemma}[definition]{Lemma}
\newtheorem{proposition}[definition]{Proposition}
\newtheorem{corollary}[definition]{Corollary}

% --- macros ---
\newcommand{\R}{\mathbb{R}}
\newcommand{\N}{\mathbb{N}}
\newcommand{\Z}{\mathbb{Z}}
\newcommand{\Vect}{\mathrm{Vect}}
\newcommand{\Multi}{\mathrm{Multi}}
\newcommand{\Inc}{\mathrm{Inc}}
\newcommand{\Attn}{\mathrm{Attn}}
\newcommand{\MLP}{\mathrm{MLP}}
\newcommand{\softmax}{\mathrm{softmax}}
\newcommand{\pool}{\mathrm{pool}}
\newcommand{\sgn}{\mathrm{sgn}}
\newcommand{\Kstruct}{K_{\mathrm{struct}}}
\newcommand{\Kafford}{K_{\mathrm{afford}}}
\newcommand{\role}[1]{\texttt{#1}}
\newcommand{\catK}{\mathcal{C}_K}

% --- title ---
\title{
Role-Labelled Incidence Complexes for Lean Proof States%
\thanks{The project is informally referred to as \emph{Topological Proof Models}; the present paper does not invoke that branding in its claims.}
\\[0.5em]
\large A Leakage-Aware Compiler and a Role-Conditioned Attention Layer
}
\author{}
\date{}

\begin{document}

\maketitle

\begin{abstract}
Formal proof states are commonly serialized as text before being processed by neural theorem provers, but their native structure is typed, scoped, relational, and higher-order. We introduce \emph{role-labelled incidence complexes} (RLICs) as a representation of Lean proof states, in which grade-$0$ cells are atomic proof objects, grade-$1$ cells are typed binary incidences, and higher-grade cells encode either structural multi-object data (term application, binder structure) or \emph{state-only affordances} (visible rewrite opportunities, visible constructor or induction targets). Boundaries are role-labelled: each incidence carries a role distinguishing the position a face plays in its coface. In general RLICs, \emph{grade} is a modelling index rather than a topological dimension; it agrees with cellular dimension only in the proof-state-cell-complex subcase. We separate four complexes, a structural complex $\Kstruct(P)$, an affordance complex $\Kafford(P)$, a candidate-action complex $K_a(P,A)$, and a transition complex $K_t(P,a,P')$, and we state a no-leakage condition on the compiler $P \mapsto \Kafford(P)$ with an explicit computational boundary on the parser. We define a \emph{role-conditioned incidence-transport} layer used by the implemented model, with face-to-coface upward maps $\rho^{r,\uparrow}$ and coface-to-face downward maps $\rho^{r,\downarrow}$, and a same-grade map indexed by a class pair. We define a \emph{cellular-sheaf} variant as a covariant functor on an incidence category whose arrows go from face to coface (Hansen--Ghrist convention); the maps are still called \emph{restriction maps} in that convention. We are explicit that in the free-incidence case the strict assignment is a quiver representation rather than a coherence-constrained sheaf. We establish: a compositional no-leakage lemma (parser compliance reduces to generator compliance); a conditional incidence-isomorphism equivariance theorem; degenerate reductions to set self-attention and to typed node, edge message passing; a purely combinatorial higher-cell sensitivity lemma, with an explicit remark that it does \emph{not} separate the model from typed-hyperedge baselines; and a kernel-soundness observation. We describe \textbf{LTM-Tiny}, a small Lean prototype, with an evaluation protocol that includes both a \emph{weak} (typed hyperedges) and a \emph{strong} (role-aware, hyperedge-stateful) hypergraph baseline, and a symbolic affordance-count baseline that tests whether tactic-family prediction is already explained by simple state-only symbolic features. Empirical results on a filtered LeanDojo Benchmark~4 subset (\S\ref{subsec:results}) corroborate the targeted hypothesis at LTM-Tiny scale: the role-conditioned RLIC layer beats a parameter-matched strong role-aware hyperedge baseline by $+1.94$~pp top-1 ($p < 10^{-3}$) and $+2.64$~pp Task-C AUC ($p < 10^{-2}$), with the largest sliced top-1 gap ($+3.19$~pp) on rewrite-heavy proof states; ablations isolate the role labels and the learnable role-conditioned transport as the two load-bearing components. To our knowledge, this specific role-labelled incidence-complex representation, with its four-complex hierarchy, leakage-aware compiler, and role-conditioned incidence-transport layer for Lean-style proof search, has not been developed previously.
\end{abstract}

% =====================================================================
\section{Introduction}
% =====================================================================

Formal proof states are not fundamentally textual objects. A proof state contains typed variables, local hypotheses, goals, metavariables, dependencies, equalities, and structural information whose form is closer to a typed higher-order incidence object than to a one-dimensional sequence. Existing graph-based representations of higher-order logic and proof states already establish that structured neural representations are useful (\S\ref{sec:related}). Our contribution is specific: a role-labelled, graded incidence-complex representation with a four-complex hierarchy separating structural state, state-only affordances, externally supplied candidate actions, and verified transitions; a parser obeying an explicit no-leakage discipline with a computational boundary; and a role-conditioned attention layer with a clearly stated cellular-sheaf variant.

\paragraph{Main contribution.} The compiler and the four-complex hierarchy, not the neural layer. The role-conditioned attention layer adapts existing topological-deep-learning machinery. What is new is the construction
\[
P \;\longmapsto\; \Kstruct(P) \;\subseteq\; \Kafford(P)
\]
under an explicit no-leakage discipline, together with the precise separation of $\Kafford(P)$ from the candidate-action complex $K_a(P,A)$ and the transition complex $K_t(P,a,P')$.

\paragraph{Terminology.} We use \emph{role-labelled incidence complex} (RLIC) rather than \emph{cell complex} by default. The constructed object does not in general satisfy the closure axioms of a regular cell complex, and the boundary operators used in cellular-sheaf theory require additional signed-orientation data that we do not impose universally. We reserve the term \emph{proof-state cell complex} for the subclass where this additional structure holds. We use \emph{grade} (rather than \emph{dimension}) for the integer label on cells in general RLICs; \emph{grade} and \emph{dimension} coincide only in the proof-state-cell-complex subcase (Definition~\ref{def:pscc}). We similarly avoid \emph{Hodge Laplacian} and \emph{cellular coboundary} outside that subcase.

\paragraph{Contributions.}
\begin{enumerate}
\item We define \emph{role-labelled incidence complexes} with role-labelled boundary multisets, with \emph{grade} as the integer cell index distinguished from cellular dimension.
\item We define a \emph{four-complex hierarchy} $\Kstruct \subseteq \Kafford \subseteq K_a$ (with $K_t$ for training only), separating pure structural cells from state-only affordances from externally supplied candidate actions.
\item We state a \emph{no-leakage condition} (Condition~\ref{cond:noleak}) with an explicit computational boundary, and we prove a \emph{compositional no-leakage lemma} (Lemma~\ref{lem:comp-noleak}) reducing parser compliance to per-generator compliance and deterministic truncation.
\item We give parser pseudo-code (Algorithms 1--4) consistent across complexes.
\item We define an \emph{incidence category} $\catK$ with face-to-coface arrows, and a \emph{cellular sheaf} as a covariant functor $\mathcal{F}\colon \catK \to \Vect_{\R}$ giving \emph{restriction maps} $\mathcal{F}(\tau) \to \mathcal{F}(\sigma)$ for $\tau \trianglelefteq \sigma$, following the cellular-sheaf convention. In the free-incidence case the strict assignment is a quiver representation (Remark~\ref{rem:free-quiver}).
\item We define a \emph{relaxed role-conditioned incidence-transport} layer with independent maps per role-and-class signature, used in the implemented model, and we are explicit that it is \emph{not} a sheaf.
\item We establish: a compositional no-leakage lemma; a conditional incidence-isomorphism equivariance theorem with explicit assumptions; reductions to set self-attention and to typed node--edge message passing; a purely combinatorial higher-cell sensitivity lemma, with an explicit remark that the lemma does not separate the model from typed-hyperedge baselines; and a kernel-soundness observation.
\item We specify policy, value, and retrieval heads in a way that avoids tactic-label leakage through the policy input; we make the retrieval-invariance caveat explicit.
\item We propose an evaluation protocol that includes a \emph{weak} and a \emph{strong} hypergraph baseline and a \emph{symbolic affordance-count} baseline that may already account for much of the tactic-family signal.
\item We report empirical results on a filtered LeanDojo Benchmark~4 subset at LTM-Tiny scale (\S\ref{subsec:results}), confirming the targeted hypothesis: M5 beats the parameter-matched strong hyperedge baseline on top-1 and Task-C AUC with three-seed variance bars, the largest sliced gap is on rewrite-heavy proof states, ablations isolate role labels and learnable role-conditioned transport as the two load-bearing components, and all three Algorithm~\ref{sec:parser}.4 no-leakage probes pass at $100\%$.
\end{enumerate}

% =====================================================================
\section{Related Work}\label{sec:related}
% =====================================================================

\paragraph{Graph-based representations for theorem proving.}
Structured neural representations of higher-order logic formulas have been used for premise selection (FormulaNet~\cite{wang2017formulanet} represents HOL terms as graphs invariant under variable renaming) and for higher-order proof search (graph representations for HOL and theorem proving~\cite{paliwal2020graph}). For Coq, Graph2Tac~\cite{blaauwbroek2024graph2tac} learns hierarchical graph representations of mathematical concepts and proof states. Recent Lean 4 work has explored graph-neural approaches connected to atomic tactic prediction. We do not claim primacy for structural proof-state representation; we contribute the specific \emph{role-labelled, graded incidence-complex} form with a four-complex hierarchy and a leakage-aware compiler.

\paragraph{Neural theorem proving with Lean.}
LeanDojo~\cite{yang2023leandojo} provides programmatic Lean interaction, extracting proof states, tactics, and premises, and supports retrieval-augmented theorem proving via ReProver. DeepSeek-Prover-V1.5~\cite{xin2024deepseek15} couples Lean~4 with proof-assistant feedback, RL, and MCTS-like search; DeepSeek-Prover-V2~\cite{deepseekv2} adds recursive subgoal decomposition. Our work is plug-compatible with such pipelines.

\paragraph{Hypergraphs and proof DAGs.}
A graded RLIC has natural reductions to typed hypergraphs (discard grading and role labels, replace higher cells with hyperedges) and to typed graphs (discard higher cells). Lemma~\ref{lem:sens} separates our layer from grade-$1$-only baselines but does \emph{not} separate it from typed-hyperedge baselines; whether \emph{graded} cellular structure further helps over typed hyperedges is an empirical question included in \S\ref{sec:experiments} via both a weak and a strong hypergraph baseline.

\paragraph{Topological deep learning.}
RLICs sit within the lineage of topological deep learning~\cite{hajij2022tdl}. \emph{Simplicial and cellular attention}~\cite{giusti2022simplicial,ballester2024cellular} define attention over higher-grade incidence; our layer adapts this to typed role-labelled patterns. \emph{Sheaf neural networks}~\cite{bodnar2022sheaf} use cellular sheaves on graphs, with standard GNNs as the trivial-sheaf case; the Hansen--Ghrist convention~\cite{hansenghrist2019} underlies our \S\ref{sec:sheaves}. Homotopy type theory offers a natural home for proof-relevant equality via higher cells; not pursued here.

% =====================================================================
\section{The Object Hierarchy}\label{sec:hierarchy}
% =====================================================================

\subsection{Role-labelled incidence complexes}

\begin{definition}[Role-labelled incidence complex]\label{def:rlic}
Let $\mathcal{R}$ be a finite set of \emph{incidence roles}, $\mathcal{L}$ a finite ontology of \emph{cell labels}, and $\mathcal{T}$ a finite ontology of \emph{type/foundation annotations}. A \emph{role-labelled incidence complex} (RLIC) is a tuple
\[
K = (K_0 \sqcup K_1 \sqcup \cdots,\; \partial,\; \ell,\; \tau),
\]
where the cell set is finite and \emph{graded} by an integer $g(\sigma) \in \Z_{\ge 0}$; $\ell\colon \bigsqcup_g K_g \to \mathcal{L}$, $\tau\colon \bigsqcup_g K_g \to \mathcal{T}$, and the \emph{role-labelled boundary map} $\partial$ sends each $\sigma$ with $g(\sigma) > 0$ to a finite multiset
\[
\partial \sigma \;\in\; \Multi\bigl(K_{<g(\sigma)} \times \mathcal{R}\bigr).
\]
\end{definition}

\begin{remark}[Grade vs dimension]\label{rem:grade}
In a general RLIC, \emph{grade} is a modelling index, not a topological dimension. We use grade $0$ for atomic objects, grade $1$ for binary incidences, and grade $\ge 2$ for multi-object structural or affordance relations. Grade-$2$ here is therefore an ``arity level'', multi-way binding, not a topological 2-cell. In the proof-state-cell-complex subcase below, grade coincides with cellular dimension.
\end{remark}

\begin{definition}[Proof-state cell complex]\label{def:pscc}
An RLIC is a \emph{proof-state cell complex} when:
\begin{enumerate}[label=(\roman*)]
\item it admits a regular cell-complex structure in which each grade-$d$ cell is topologically $d$-dimensional;
\item each cell carries an orientation;
\item the boundary multiset $\partial \sigma$ for a $d$-cell $\sigma$ refines into a distinguished \emph{codimension-one} sub-multiset $\partial^{(1)} \sigma \subseteq \partial \sigma$ consisting of pairs $(\tau, r)$ with $g(\tau) = d - 1$ together with signs derived from the orientation;
\item the signed codimension-one boundary $d_d\colon C_d \to C_{d-1}$ extracted from $\partial^{(1)}$ satisfies $d_{d-1} \circ d_d = 0$.
\end{enumerate}
Lower-codimension entries in $\partial \sigma$ (those with $g(\tau) < d - 1$) are treated in this subcase as auxiliary \emph{role-annotated incidences}, not as direct cellular boundary terms.
\end{definition}

We do \emph{not} assume RLICs are proof-state cell complexes by default. Any statement that depends on signed boundary operators, chain or cochain complexes, $\delta^2 = 0$, or Hodge/sheaf Laplacians is made only in this subcase.

\subsection{Roles, orientation, and adjacency}

\begin{definition}[Role set]\label{def:roles}
The role set used in this paper is
\begin{align*}
\mathcal{R} = \{ & \role{source}, \role{target}, \role{function}, \role{argument}_k, \role{goal}, \role{equality}, \role{occurrence},\\
& \role{binder}, \role{bound}, \role{premise}, \role{conclusion}, \role{term}, \role{type}, \role{proof},\\
& \role{proposition}, \role{inductive}, \role{base\_case}, \role{step\_case}, \role{witness} \},
\end{align*}
with $\role{argument}_k$ indexed by $k \in \N$.
\end{definition}

Each cell \emph{class}, determined by $(\ell(\sigma), \tau(\sigma), g(\sigma))$, fixes the role signature of its boundary.

\paragraph{Adjacency.} Two cells of the same grade are \emph{adjacent} if they share at least one boundary entry; they are \emph{role-pair-adjacent} if they share a face $\eta$ with role pair $(r, r')$, where $(\eta, r) \in \partial \sigma$ and $(\eta, r') \in \partial \sigma'$.

\subsection{Grade-0, grade-1, and higher-grade cells}

\paragraph{Grade-0 cells} are atomic proof objects: variables, constants, local hypotheses, the goal proposition, type and universe expressions, terms, selected subterms, metavariables, and referenced theorem or lemma names. The boundary is empty.

\paragraph{Grade-1 cells} have boundary multisets of size two, recording typed binary incidences: \role{has-type}, \role{occurs-in}, \role{proves}, \role{binds}, etc., each with the role pair appropriate to the relation.

\paragraph{Higher-grade cells} come in two distinct categories:
\begin{itemize}
\item \emph{Structural higher-grade cells}: cells that encode multi-way \emph{state-side} structure, for instance, a grade-$2$ application cell with roles $(\role{function}, \role{argument}_1, \ldots, \role{argument}_n, \role{conclusion})$ representing the application $f\,a_1\,\cdots\,a_n$ in the proof state itself. These cells exist because the \emph{state contains} an application, not because of any potential next move.
\item \emph{Affordance higher-grade cells}: cells recording \emph{state-only} candidate inference opportunities, for instance, a \role{rewrite-opp} grade-$2$ cell when an equality in the local context has at least one occurrence of its LHS or RHS in the goal. These cells are emitted by fixed, $P$-independent generators (\S\ref{sec:noleak}).
\end{itemize}

\subsection{The four complexes}

\begin{definition}[Four complexes]\label{def:four}
Given $(P, A, a, P')$ with $a \in A$ and $P'$ the successor of $P$ under $a$:
\begin{itemize}
\item $\Kstruct(P)$ contains only grade-$0$ cells, grade-$1$ cells, and structural higher-grade cells, all determined by the parse tree, typing, and local context of $P$.
\item $\Kafford(P) \supseteq \Kstruct(P)$ additionally contains affordance higher-grade cells, each emitted by a fixed affordance generator from data inside $P$.
\item $K_a(P, A) \supseteq \Kafford(P)$ additionally contains one \emph{candidate-action cell} per element of $A$, with role-labelled boundary lying entirely in $\Kafford(P)$.
\item $K_t(P, a, P')$ is the transition complex obtained after the kernel accepts $a$ and produces $P'$; it is used only for training transition models and is never available to the policy, value, or retrieval heads at prediction time.
\end{itemize}
\end{definition}

Value and retrieval heads consume $\Kafford(P)$. The policy head consumes $\Kafford(P)$ in Mode~1 (tactic-family classification) and $K_a(P, A)$ in Mode~2 (candidate-instance scoring); see \S\ref{sec:heads}.

\subsection{The no-leakage condition}\label{sec:noleak}

\begin{condition}[No-leakage with computational boundary]\label{cond:noleak}
A parser $P \mapsto \Kafford(P)$ satisfies the \emph{no-leakage condition} if:
\begin{enumerate}[label=(\roman*)]
\item \textbf{Determinism.} It is a deterministic function of $P$ and a fixed parser configuration $\Pi$ (label/role ontology, depth bounds, affordance generators, truncation policy).
\item \textbf{Computational boundary.} Allowed queries on $P$:
\begin{enumerate}[label=(\alph*)]
\item syntactic inspection of the parse tree, names, and binders of the local context, hypotheses, and goal;
\item Lean kernel typing queries at a fixed transparency level $\Theta$ specified in $\Pi$;
\item syntactic-only unification (first-order matching) for affordance pattern detection;
\item typeclass instance lookup if explicitly enabled in $\Pi$, at bounded depth.
\end{enumerate}
Forbidden: tactic execution, elaboration of candidate tactic scripts, kernel reductions outside $\Theta$, successor metavariable states, and any data not derivable from $P$ under (a) (d).
\item \textbf{Independence of next move.} $\Kafford(P)$ does not depend on the next human tactic, the next state $P'$, the lemma applied next, or any retrieval shortlist whose construction violates (i) (iii).
\end{enumerate}
The candidate set $A$ used in $K_a(P, A)$ is generated either by a fixed \emph{family-universal affordance generator} (emitting the same set of structural candidate slots regardless of the next tactic) or by a retrieval shortlist itself satisfying (i)--(iii).
\end{condition}

\subsection{Parser algorithms}\label{sec:parser}

\begin{small}
\begin{verbatim}
Algorithm 1: BuildStruct(P, Pi)
  K_0 <- empty; K_1 <- empty; K_{>=2} <- empty

  // grade-0 cells
  for each variable v in Gamma:               emit sigma_v
  for each (h : T) in Gamma:                  emit sigma_h and sigma_T
  emit sigma_g for the goal proposition
  for each named constant c referenced:       emit sigma_c
  for each subterm s up to depth bound D:     if novel, emit sigma_s

  // grade-1 cells (each with role-labelled boundary multiset)
  for each typing (t : T) visible:
    emit "has-type",  d = {(sigma_t, term), (sigma_T, type)}
  for each occurrence of v in expression e:
    emit "occurs-in", d = {(sigma_v, occurrence), (sigma_e, prop|term)}
  for each (h : T) in Gamma:
    emit "proves",    d = {(sigma_h, proof), (sigma_T, proposition)}
  for each binder beta with body b:
    emit "binds",     d = {(sigma_beta, binder), (sigma_b, bound)}

  // structural higher-grade cells (grade >= 2)
  for each application f a_1 ... a_n:
    emit grade-2 "app-(n)",
      d = {(sigma_f, function), (sigma_{a_1}, argument_1), ...,
           (sigma_{a_n}, argument_n), (sigma_result, conclusion)}
  for each equality literal a = b:
    emit grade-2 "eq-lit",
      d = {(sigma_a, source), (sigma_b, target), (sigma_Eq, function)}
  // further structural families per Pi

  return K_struct(P)
\end{verbatim}
\end{small}

\begin{small}
\begin{verbatim}
Algorithm 2: BuildAfford(P, Pi)
  K <- BuildStruct(P, Pi)
  // Each generator G_i is a deterministic function of (P, Pi) using only
  // queries allowed by Condition 3.6(ii).

  for each equality hypothesis h : a = b:
    for each occurrence of a (resp. b) in Gamma U {g}:
      emit "rewrite-opp",  d = {(sigma_h, equality), (sigma_occ, occurrence)}

  if head of g is constructive (and, or, exists, iff):
    emit "constructor-opp",
      d = {(sigma_g, goal), (sigma_{f_1}, argument_1), ..., (sigma_{f_k}, argument_k)}

  for each inductive variable v whose type appears in g:
    emit "induction-opp", d = {(sigma_v, inductive), (sigma_g, goal)}

  for each (h, g) pair where head(type(h)) matches head(g)
       via first-order matching at Theta:
    emit "apply-opp", d = {(sigma_h, premise), (sigma_g, conclusion), ...}

  return K
\end{verbatim}
\end{small}

\begin{small}
\begin{verbatim}
Algorithm 3: BuildCandidateAction(P, A, Pi)
  K <- BuildAfford(P, Pi)
  for each a in A:
    B_a <- MatchBoundary(a, K)         // draws faces only from K_afford(P)
    emit c_a with signature(a), d(c_a) = B_a
  return K
\end{verbatim}
\end{small}

\begin{small}
\begin{verbatim}
Algorithm 4: NoLeakageProbes(parser, dataset)
  // Necessary, not sufficient, evidence for no leakage.

  // Probe 1: Tactic-swap.
  for each (P, a_human, P'):
    K1, K2 <- parser.BuildAfford(P), parser.BuildAfford(P)
    assert K1 == K2 up to canonical ordering
    for each alternative a' != a_human:
      K3 <- parser.BuildAfford(P)
      assert K1 == K3                  // a' must not influence K_afford

  // Probe 2: Permuted-context (alpha-renaming).
  for each P:
    sigma <- random label-class-preserving renaming of bound variables
    K1, K2 <- parser.BuildAfford(P), parser.BuildAfford(sigma(P))
    assert K2 ~= sigma_*(K1)

  // Probe 3: Held-out lemma in retrieval shortlist.
  for each (P, A, a_human):
    A_minus <- A \ {lemma(a_human)}
    Ka  <- parser.BuildCandidateAction(P, A)
    Ka' <- parser.BuildCandidateAction(P, A_minus)
    assert K_afford(Ka) == K_afford(Ka')
    assert Ka, Ka' differ only by candidate-action cells tied to lemma(a_human)

  return pass-rate, failing examples
\end{verbatim}
\end{small}

\subsection{Worked example}

For \texttt{x y : Nat; h : x = y \(\vdash\) y = x}:

$\Kstruct(P)$: grade-$0$ cells for \texttt{x}, \texttt{y}, \texttt{Nat}, \texttt{h}, \texttt{x = y}, \texttt{y = x}, \texttt{Eq}. Grade-$1$ cells for $\role{has-type}(\cdot, \texttt{Nat})$, $\role{proves}(h, \texttt{x = y})$, and \role{occurs-in} cells. Two grade-$2$ \role{eq-lit} cells for \texttt{x = y} and \texttt{y = x}, each with $\{(\text{lhs}, \role{source}), (\text{rhs}, \role{target}), (\sigma_{\mathrm{Eq}}, \role{function})\}$.

$\Kafford(P) \setminus \Kstruct(P)$: two \role{rewrite-opp} grade-$2$ cells (for the occurrences of \texttt{x} and \texttt{y} in the goal); one \role{apply-opp} grade-$2$ cell with $\{(\sigma_h, \role{premise}), (\sigma_{y=x}, \role{conclusion})\}$.

$K_a(P,A) \setminus \Kafford(P)$: one candidate-action cell per element of $A$ (\texttt{rw [h]}, \texttt{rw [\(\leftarrow\) h]}, \texttt{exact h.symm}, \texttt{apply Eq.symm}, \texttt{rfl}, \ldots), boundaries lying in $\Kafford$.

What is absent: no \texttt{y = y}, no post-rewrite expression, no cell whose existence depends on the next tactic or successor state.

% =====================================================================
\section{Incidence Sheaves and Transport}\label{sec:sheaves}
% =====================================================================

\subsection{The incidence category}

\begin{definition}[Incidence category $\catK$]\label{def:catK}
Given an RLIC $K$, $\catK$ has objects = cells; generating morphisms $\iota^r_{\tau \to \sigma}\colon \tau \to \sigma$ for each $(\tau, r) \in \partial \sigma$ (face to coface, role-labelled); composition = free composition of generators, optionally quotiented by parser-supplied path identifications.
\end{definition}

\subsection{Cellular sheaves on $\catK$}

\begin{definition}[Cellular sheaf]\label{def:sheaf}
A \emph{cellular sheaf} on $K$ is a covariant functor
\[
\mathcal{F}\colon \catK \longrightarrow \Vect_{\R},
\]
assigning a finite-dimensional stalk $\mathcal{F}(\sigma)$ and, to each role-labelled incidence, a linear map
\[
\rho^r_{\tau \trianglelefteq \sigma}\;\colon\; \mathcal{F}(\tau) \to \mathcal{F}(\sigma),
\]
satisfying functoriality on every composable path in $\catK$.
\end{definition}

\paragraph{Convention.} Although these maps go from face to coface, we follow the cellular-sheaf convention of calling them \emph{restriction maps}; this matches Hansen--Ghrist~\cite{hansenghrist2019} and the cellular-sheaf literature, and differs from the ordinary topological sheaf convention where restriction goes from larger open sets to smaller ones.

\begin{remark}[Triviality in the free-incidence case]\label{rem:free-quiver}
If $\catK$ is freely generated, any assignment of vector spaces to objects and linear maps to generating morphisms uniquely determines a functor; the functoriality is automatic. Nontrivial coherence appears only when the parser identifies distinct composable paths between the same pair of cells. The free case is the default; nontrivial sheaf structure is parser-imposed.
\end{remark}

\subsection{Coboundary and incidence Dirichlet energy (restricted setting)}

\begin{definition}[Coboundary and Dirichlet energy in the proof-state-cell-complex subcase]\label{def:dirichlet}
When $K$ is a proof-state cell complex (Definition~\ref{def:pscc}) and $\mathcal{F}$ is a cellular sheaf in which restriction maps are compatible with the chosen orientation in the standard cellular-sheaf sense (i.e., signs on codimension-one incidences match the signed boundary operator $d_d$ extracted from $\partial^{(1)}$), the cellular-sheaf coboundary $\delta^d\colon C^d(K, \mathcal{F}) \to C^{d+1}(K, \mathcal{F})$ on cochains is given by
\[
(\delta^d c)(\sigma) \;=\; \sum_{(\tau, r) \in \partial^{(1)} \sigma} \sgn(\tau, \sigma)\,\rho^r_{\tau \trianglelefteq \sigma}\bigl(c(\tau)\bigr),
\]
with $\sgn(\tau, \sigma)$ the orientation sign on the codimension-one incidence. The sheaf Laplacian is $L = \delta^* \delta$ and the \emph{incidence Dirichlet energy} of a $0$-cochain $x$ is $\langle x, Lx \rangle$. Role labels on incidences of codimension $>1$ do not enter $\delta^d$; they remain auxiliary annotations.
\end{definition}

We invoke Definition~\ref{def:dirichlet} \emph{only} in this subcase. The implemented model does not.

\subsection{Role-conditioned incidence transport (the implemented model)}

\begin{definition}[Relaxed role-conditioned incidence transport]\label{def:transport}
A \emph{relaxed incidence transport} on an RLIC $K$ assigns to each role-labelled incidence $(\tau, r) \in \partial \sigma$ two independently parameterised maps:
\begin{align*}
\rho^{r,\uparrow}_{\tau \to \sigma} &\colon \R^{d_\tau} \to \R^{d_\sigma} \quad\text{(face-to-coface, ``upward''),}\\
\rho^{r,\downarrow}_{\sigma \to \tau} &\colon \R^{d_\sigma} \to \R^{d_\tau} \quad\text{(coface-to-face, ``downward''),}
\end{align*}
each realised as a small MLP whose parameters are shared across all incidences of the same role-and-class signature, and \emph{not} required to satisfy any functoriality, adjointness, or composition law.
\end{definition}

\begin{remark}\label{rem:not-sheaf}
The relaxed transport is \textbf{not a cellular sheaf}. It does not assemble into a functor; it carries an independently parameterised downward map; and it discards orientation/sign data. Statements that depend on sheaf structure (e.g.\ Definition~\ref{def:dirichlet}) apply only in the strict setting.
\end{remark}

\subsection{Incidence transport residual (relaxed)}

For the relaxed model:
\[
R(x) \;=\; \sum_{(\sigma, \tau, r)\in \Inc(K)} \Bigl\lVert \rho^{r,\uparrow}_{\tau \to \sigma}(x_\tau) - x_\sigma \Bigr\rVert^2,
\]
a learnable structural-coherence signal usable as an auxiliary loss. It does not prove correctness; the Lean kernel does.

% =====================================================================
\section{Architecture}\label{sec:arch}
% =====================================================================

\subsection{Cell hidden states}

Hidden states $h_\sigma \in \R^{d_C}$ per cell-class signature $C$ determined by $(\ell(\sigma), \tau(\sigma), g(\sigma))$. By default all $d_C = d$; cross-class transport maps then become same-dimensional linear maps but still differ by role.

\subsection{Neighbourhoods}

For each cell $\sigma$:
\begin{itemize}
\item $\mathcal{N}_{\downarrow}(\sigma) = \{(\tau, r) : (\tau, r) \in \partial \sigma\}$ \emph{role-labelled boundary} (faces of $\sigma$).
\item $\mathcal{N}_{\uparrow}(\sigma) = \{(\omega, r) : (\sigma, r) \in \partial \omega\}$ \emph{role-labelled coboundary} (cofaces of $\sigma$).
\item $\mathcal{N}_{\leftrightarrow}(\sigma) = \{(\sigma', \eta, r, r') : \sigma' \neq \sigma,\, g(\sigma') = g(\sigma),\, (\eta, r) \in \partial \sigma,\, (\eta, r') \in \partial \sigma'\}$ \emph{same-grade adjacency}, indexed by the shared face $\eta$ and the role pair $(r, r')$.
\end{itemize}
Repeated entries are treated as distinct neighbours.

\subsection{Role-conditioned attention (transport directions made explicit)}\label{sec:attention}

We are explicit about which transport map serves which neighbourhood, since both directions of $\rho$ exist (Definition~\ref{def:transport}):
\begin{itemize}
\item For $(\tau, r) \in \mathcal{N}_{\downarrow}(\sigma)$ $\tau$ is a face of $\sigma$ the message into $\sigma$ uses the \emph{upward} map $\rho^{r,\uparrow}_{\tau \to \sigma}(h_\tau)$ (transporting face $\to$ coface).
\item For $(\omega, r) \in \mathcal{N}_{\uparrow}(\sigma)$ $\omega$ is a coface of $\sigma$ the message into $\sigma$ uses the \emph{downward} map $\rho^{r,\downarrow}_{\omega \to \sigma}(h_\omega)$ (transporting coface $\to$ face).
\item For $(\sigma', \eta, r, r') \in \mathcal{N}_{\leftrightarrow}(\sigma)$  $\sigma'$ is a same-grade neighbour sharing the face $\eta$ the message into $\sigma$ uses a same-grade map
\[
\rho^{(r, r', \ell(\eta), \tau(\eta), C(\sigma'), C(\sigma))}_{\sigma' \to \sigma}(h_{\sigma'}),
\]
i.e., one MLP per signature consisting of the two roles, the shared-face label and type, and the source/target cell classes.
\end{itemize}

The packaged attention formula per neighbourhood direction $\bullet \in \{\downarrow, \uparrow, \leftrightarrow\}$:
\[
q_\sigma = W_Q^{\bullet} h_\sigma,\quad k^{\bullet}_{\nu \to \sigma} = W_K^{\bullet} \, m^{\bullet}_{\nu \to \sigma},\quad v^{\bullet}_{\nu \to \sigma} = W_V^{\bullet}\, m^{\bullet}_{\nu \to \sigma},
\]
where $m^{\bullet}_{\nu \to \sigma}$ is the appropriate transported message defined above, $\nu$ ranges over the elements of $\mathcal{N}_{\bullet}(\sigma)$, and
\[
\alpha^{\bullet}_{\sigma\nu} = \softmax_\nu\!\left(\frac{q_\sigma^\top k^{\bullet}_{\nu \to \sigma}}{\sqrt{d}}\right),\qquad \Attn_{\bullet}(\sigma) = \sum_\nu \alpha^{\bullet}_{\sigma\nu}\, v^{\bullet}_{\nu \to \sigma}.
\]
The $W_K, W_V$ matrices may additionally be role-indexed within each $\bullet$.

\subsection{Layer update}

\[
h_\sigma^{\ell+1} = \MLP\bigl(h_\sigma^\ell,\, \Attn_{\downarrow}(\sigma),\, \Attn_{\uparrow}(\sigma),\, \Attn_{\leftrightarrow}(\sigma)\bigr),
\]
with residual connections and pre-layernorm.

\subsection{Labels, embeddings, and truncation (equivariance-critical)}\label{sec:labels}

To make Theorem~\ref{thm:equiv} hold:
\begin{itemize}
\item \textbf{Class-shared embeddings.} Cell embeddings depend on $(\ell(\sigma), \tau(\sigma), g(\sigma))$, not cell identity.
\item \textbf{Canonical truncation.} When budgets are exceeded, cells are discarded by a canonical label-and-incidence-only policy so that isomorphic complexes truncate to isomorphic complexes.
\item \textbf{Equivariant positional features.} Sign- and basis-invariant aggregations only; naive eigenvector entries not used.
\item \textbf{Invariant pooling.} Sum or class-canonical attention pooling.
\end{itemize}

\paragraph{Label preservation policy.} $\alpha$-equivalence requires care.

\begin{center}
\begin{tabular}{p{0.40\linewidth} p{0.18\linewidth} p{0.34\linewidth}}
\toprule
\textbf{Entity} & \textbf{Raw name preserved?} & \textbf{Reason} \\
\midrule
bound variable & no & $\alpha$-equivalence \\
local hypothesis name (\texttt{h}, \texttt{h\_1}, \texttt{this}) & no by default; optional via $\Pi$ & avoid name leakage from style \\
global theorem / lemma constant & yes & semantic premise identity \\
inductive type / constructor names & yes & semantic identity \\
namespace / module path & configurable; off by default & dataset-leakage risk \\
generated metavariable name & no & artefact \\
universe-level variable & no & artefact \\
\bottomrule
\end{tabular}
\end{center}

When a raw name is preserved, the embedding goes through a fixed vocabulary lookup; when normalised, through a class embedding plus a $\Pi$-dependent canonical key (e.g. de~Bruijn index for bound variables).

\subsection{Heads (input choice prevents Mode-1 label leakage)}\label{sec:heads}

\paragraph{Mode 1 (tactic-family classification).} Input: $\Kafford(P)$, \emph{not} $K_a(P, A)$. Output: categorical distribution over tactic families. The policy pools class-equivariantly:
\[
\pi(\text{family} \mid \Kafford(P)) = \softmax\bigl(W_\pi \cdot \pool_{\mathrm{class}}(h^L)\bigr).
\]
\emph{No candidate-action cells appear in the input:} the model cannot read off the answer from family-labelled cells.

\paragraph{Mode 2 (candidate-instance scoring).} Input: $K_a(P, A)$, with each candidate-action cell carrying a \emph{signature} (structural data of $a$) that does not encode a separate tactic-family target label. The candidate generator must be family-universal and tactic-independent. We do not commit to one in this paper.

\paragraph{Value head.} Pools $\Kafford(P)$ to a scalar in $[0, 1]$.

\paragraph{Retrieval head.} $\mathrm{score}(\Kafford(P), t_i) = \phi(\Kafford(P))^\top \psi(t_i)$ for a premise $t_i$. We discuss the invariance caveat in Theorem~\ref{thm:equiv}.

% =====================================================================
\section{Formal Properties}\label{sec:theory}
% =====================================================================

\begin{definition}[Proof-complex isomorphism]\label{def:iso}
$\varphi\colon K \to K'$ over the same ontology is a \emph{proof-complex isomorphism} if it is a bijection of cells with (a) $g \circ \varphi = g$, (b) $(\ell', \tau') \circ \varphi = (\ell, \tau)$, and (c) the multiset equality $\partial' \varphi(\sigma) = \{(\varphi(\tau), r) : (\tau, r) \in \partial \sigma\}$ for all $\sigma$. Labels are modulo the preservation policy of \S\ref{sec:labels}.
\end{definition}

\begin{lemma}[Compositional no-leakage]\label{lem:comp-noleak}
Let $\Pi$ have affordance generators $G_1, \ldots, G_m$, structural generator $G_{\mathrm{struct}}$, and truncation policy $T$. Suppose:
\begin{enumerate}[label=(\roman*)]
\item $G_{\mathrm{struct}}$ and each $G_i$ are deterministic functions of $(P, \Pi)$ using only queries allowed by Condition~\ref{cond:noleak}(ii);
\item $T$ is deterministic in $(\Kstruct(P), \Pi)$ and uses no information about the next move;
\item the candidate set $A$ is generated by a family-universal procedure, or by a retrieval shortlist itself obeying (i)--(ii).
\end{enumerate}
Then the composed parser (Algorithms~1--3 followed by $T$) satisfies Condition~\ref{cond:noleak}.
\end{lemma}

\begin{proof}
(i) Deterministic functions compose to a deterministic function, so the composed parser is deterministic in $(P, \Pi)$. (ii) Each component uses only allowed queries; the union of allowed queries used by any subset of components remains allowed; therefore the composed parser respects the computational boundary. (iii) None of $G_{\mathrm{struct}}, G_1, \ldots, G_m$, or $T$ uses the next move; $A$'s construction by hypothesis does not; therefore $\Kafford(P)$ and its truncated form are independent of the next move. The three conditions of \ref{cond:noleak} are satisfied.
\end{proof}

This is the working contract: certify each generator and the truncation policy, get parser compliance for free.

\begin{theorem}[Conditional incidence-isomorphism equivariance]\label{thm:equiv}
Suppose the architecture of \S\ref{sec:arch} is instantiated under:
\begin{enumerate}[label=(\roman*)]
\item initialization assigns embeddings depending only on $(\ell(\sigma), \tau(\sigma), g(\sigma))$;
\item labels follow the preservation policy of \S\ref{sec:labels};
\item all parameter weights are indexed by class-and-role signatures, not by cell identity;
\item truncation is canonical (\S\ref{sec:labels});
\item positional features, if used, are constructed equivariantly;
\item pooling is class-canonical.
\end{enumerate}
Then for every proof-complex isomorphism $\varphi\colon K \to K'$ and every layer $\ell$,
\[
h^\ell_{K'}(\varphi(\sigma)) \;=\; h^\ell_{K}(\sigma) \qquad \text{for all } \sigma \in K.
\]
Consequently the policy output (Mode~1 or Mode~2) and the value output are invariant under $\varphi$.
\end{theorem}

\paragraph{Retrieval caveat.} The retrieval score $\phi(K)^\top \psi(t_i)$ is invariant under $\varphi$ for a fixed candidate premise set with an isomorphism-invariant state encoder $\phi$. Invariance of the \emph{ranking} over a candidate set holds when $\psi$ is computed independently of $K$. If $\psi$ is itself a function of $K$ (e.g. cross-encoder retrieval), an additional consistency assumption on $\psi$ is required.

\begin{proof}
Under (i)--(iii) every parameter and every initial hidden state is a function of class-and-role data and of the role-labelled boundary structure, both preserved by $\varphi$. The \S\ref{sec:arch} update on $K'$ at $\varphi(\sigma)$ is therefore evaluated on the same arguments as the update on $K$ at $\sigma$. Under (iv) truncation commutes with $\varphi$. Under (v) positional features are isomorphism-invariant. Under (vi) pooling commutes with relabelling. Induct on layers. The retrieval caveat follows from inspection of $\phi(K)^\top \psi(t_i)$ given the conditions on $\psi$.
\end{proof}

\begin{proposition}[Set self-attention reduction]\label{prop:set}
If $K = K_0$ with same-grade adjacency the complete graph on $K_0$, the \S\ref{sec:arch} update reduces to permutation-equivariant set self-attention on $K_0$. Standard Transformer self-attention is recovered up to an \emph{externally supplied} total order on $K_0$ with corresponding positional features; in that case proof-complex isomorphism invariance is restricted to the order-preserving subgroup.
\end{proposition}

\begin{proposition}[Typed node--edge MP reduction]\label{prop:nodeedge}
If $K = K_0 \sqcup K_1$, the \S\ref{sec:arch} update reduces to a typed node--edge message-passing scheme on the incidence graph with edges carrying hidden states and standard node--edge incidence, additionally role-conditioned. This is a reduction to a specific sub-class of MP-GNNs, not to all of them.
\end{proposition}

\begin{lemma}[Higher-grade sensitivity, purely combinatorial]\label{lem:sens}
There exist RLICs $K, K'$ over the same ontology such that:
\begin{enumerate}[label=(\alph*)]
\item the labelled, role-labelled grade-$1$ skeleta of $K$ and $K'$ are isomorphic;
\item $K$ contains a grade-$2$ cell $c$ with role-labelled boundary $B \subseteq K_0 \cup K_1$, while $K'$ contains no grade-$2$ cell with that role-labelled boundary;
\item there exists a parameter setting of the \S\ref{sec:arch} layer (under the assumptions of Theorem~\ref{thm:equiv}) under which the outputs on $K$ and $K'$ are not isomorphic.
\end{enumerate}
Consequently, no model factoring through the labelled, role-labelled grade-$1$ skeleton can match the layer's outputs on this pair.
\end{lemma}

\begin{proof}
Pick $u, v \in K_0$ in some class admitting a \role{rewrite-opp}-like grade-$2$ cell. Let $K'$ have only $K_0 \sqcup K_1$ with no grade-$\ge 2$ cells; let $K$ extend $K'$ by a single grade-$2$ cell $c$ with $\partial c = \{(u, r_1), (v, r_2)\}$. The grade-$1$ skeleta agree. Choose $\rho^{r_1, \downarrow}_{c \to u}$ nonzero, all other relevant weights zero in same-grade and coboundary-to-$u$ directions. In $K$, by \S\ref{sec:attention} the update of $h_u$ is driven by the coboundary message $\rho^{r_1, \downarrow}_{c \to u}(h_c) \neq 0$; in $K'$ this contribution is absent (no such coface $c$). For generic remaining weights the final-layer pooling preserves the difference.
\end{proof}

\begin{remark}\label{rem:hyperedge}
Lemma~\ref{lem:sens} separates the \S\ref{sec:arch} layer from any model factoring through the grade-$1$ skeleton. It does \textbf{not} separate it from a typed-hyperedge model that encodes $c$ as a typed hyperedge over its boundary entries. Whether \emph{graded} cellular structure plus boundary/coboundary/same-grade attention beats a strong role-aware hyperedge baseline is an empirical question (\S\ref{sec:experiments}).
\end{remark}

\begin{remark}[Proof-state instantiation]\label{rem:instantiation}
Affordance generators can produce $P, P'$ with isomorphic grade-$1$ skeleta but differing affordance cells. The combinatorial lemma then applies. We do not claim this happens in any given dataset without empirical evidence.
\end{remark}

\begin{observation}[Kernel soundness]\label{obs:kernel}
If every proposed transition is checked by the Lean kernel before acceptance, the composed pipeline (LTM proposer $\circ$ kernel checker) cannot accept an invalid proof.
\end{observation}

A discipline statement, not a theorem.

% =====================================================================
\section{LTM-Tiny: Prototype and Evaluation Protocol}\label{sec:experiments}
% =====================================================================

\subsection{Dataset}

A filtered subset of LeanDojo's \texttt{mathlib} extractions; algebra, logic, and basic numerical libraries; bounded proof length and cell budgets.

\subsection{Tasks (Mode 1 throughout)}

\begin{itemize}
\item \textbf{Task A tactic-family prediction.} Input: $\Kafford(P)$ (not $K_a$). Metrics: top-1, top-5, cross-entropy.
\item \textbf{Task B premise retrieval.} Input: $\Kafford(P)$. Metrics: Recall@10, Recall@50, MRR.
\item \textbf{Task C value prediction.} Input: $\Kafford(P)$. Metrics: AUC, accuracy, calibration error.
\end{itemize}

\subsection{Baselines}\label{sec:baselines}

\paragraph{Symbolic affordance-count baseline.}
A feature vector of per-class affordance-cell counts plus basic syntactic features (goal head, number of hypotheses, presence of equalities, etc.), fed to logistic regression and to a small MLP. This baseline tests whether tactic-family prediction is already explained by symbolic features extracted by the affordance generators  in which case the neural architecture's contribution lies elsewhere (retrieval, value, or harder structural cases) rather than in basic family classification.

\paragraph{Text baseline.}
A text Transformer over the pretty-printed proof state.

\paragraph{Weak hypergraph baseline.}
Each higher-grade cell in $\Kafford$ becomes a typed hyperedge; no role conditioning, no hyperedge hidden states. This is the ``deliberately weak'' baseline that Lemma~\ref{lem:sens} \emph{does} separate against.

\paragraph{Strong role-aware hypergraph baseline.}
Each higher-grade cell becomes a typed hyperedge with role-labelled incidence positions. The model maintains hidden states for both nodes (grade-$0$ and grade-$1$ cells) and hyperedges, with bipartite node--hyperedge message passing parameterised by one MLP per role-and-class signature, matching our model's role-conditioned transport in expressive power for node--hyperedge attention. No grading, no boundary/coboundary/same-grade attention, no cellular-sheaf variant. This baseline isolates whether \emph{graded cellular structure with directional attention} helps beyond a strong role-aware hyperedge representation. We match parameter count.

\paragraph{Role-typed GNN.}
Over the grade-$0$/grade-$1$ incidence graph only; no higher cells.

\paragraph{Structural-only RLIC.}
$\Kstruct$ used instead of $\Kafford$, with our full layer.

\paragraph{Full RLIC (the model).}
$\Kafford$ with the relaxed role-conditioned transport.

\paragraph{Ablations.}
(a) collapse all roles into one;
(b) freeze $\rho^{r,\uparrow}, \rho^{r,\downarrow}$ to identity;
(c) drop incidence-spectral PEs;
(d) drop $\mathcal{N}_\leftrightarrow$;
(e) strict cellular-sheaf variant on parser configurations with imposed path identifications.

\subsection{LTM-Tiny configuration}

$d = 256$; $L = 4$ layers; $H = 4$ heads per neighbourhood; per-signature 2-layer transport MLPs with class-and-role weight sharing; budgets $(|K_0|, |K_1|, |K_{\ge 2}|) = (256, 512, 64)$; canonical truncation per \S\ref{sec:labels}.

\subsection{Sliced evaluation}

All metrics sliced by structural category rewrite-heavy, multi-argument application, constructor-split, induction, residual to test the targeted hypothesis directly.

\subsection{No-leakage probes}

Algorithm~4 is run on the parser before any training. Per Lemma~\ref{lem:comp-noleak}, parser compliance can be certified compositionally; the probes are an empirical drift check, \emph{necessary but not sufficient}. Probe pass-rate is reported alongside task results.

\subsection{Targeted hypothesis}

\begin{quote}
Adding higher-grade cells with role-conditioned transport improves performance on structurally complex proof states relative to text-only, plain-graph, \emph{weak} hyperedge, \emph{strong role-aware} hyperedge, and \emph{symbolic affordance-count} baselines, with little or no benefit on structurally trivial states; ablations that collapse roles or freeze transport degrade on rewrite- and application-heavy slices in particular.
\end{quote}

This is the strong-baseline version of the hypothesis. If the symbolic affordance-count baseline already matches the model on Task~A, the model's contribution lies in Tasks~B and~C and in the structurally complex slices of Task~A, not in basic family prediction.

\medskip
\noindent (\emph{Outcome at LTM-Tiny scale; see \S\ref{subsec:results}.) The \texttt{rewrite\_heavy} prediction is corroborated; the ``little or no benefit on structurally trivial states'' part is not, as M5 also wins the \texttt{residual} slice by $+2.35$~pp.}

\subsection{Empirical results}\label{subsec:results}

We instantiate LTM-Tiny exactly as in \S\ref{sec:experiments} and report results
on a filtered LeanDojo Benchmark~4 subset (random split, $7$ file-path
prefixes: \texttt{Mathlib/Algebra/}, \texttt{Mathlib/Logic/},
\texttt{Mathlib/Data/\{Nat,Int,Rat,Real\}/}, \texttt{Mathlib/Order/}; proof
length $\le 30$, goal text $\le 1024$ chars; synthesised Task~C negative
states at $30\%$ rate via cross-theorem goal swap). After filtering this
yields $45{,}450$ train, $10{,}414$ validation, and $10{,}406$ test
ProofStates; the global premise pool from \texttt{corpus.jsonl} has size
$16{,}384$. State parsing uses \texttt{lean\_dojo.parse\_goals} for the
hypothesis/goal split and a recursive expression parser for term structure.
M3 is parameter-matched to M5 at $d_{\mathrm{M3}} = 448$ ($35.81$~M
parameters vs M5's $36.04$~M, $\Delta = +0.6\%$). Training: $80{,}000$
optimisation steps, batch $32$ (grad-accum $1$), AdamW
$\eta = 3 \times 10^{-4}$ with cosine schedule and $800$ warmup steps, BF16
mixed precision; Task~B is trained with $512$ random global negatives per
step in addition to in-batch negatives. Hardware: single NVIDIA GeForce
RTX~4060~Ti 16~GB; total wall-clock for the full sweep (Phase~1 + two
additional seeds on M3 and M5 + Phase~3 ablations, twelve subprocess runs
in all) was $17.54$~hours.

\subsubsection*{Algorithm~\protect\ref{sec:parser}.4 no-leakage probes (Lemma~\protect\ref{lem:comp-noleak})}

Probes were run on a $1{,}000$-state validation sample after parser
construction and before any training:

\begin{center}
\begin{tabular}{lcc}
\toprule
Probe & Pass / total & Pass-rate \\
\midrule
Tactic-swap (Probe~1) & $1000 / 1000$ & $100.00\%$ \\
Permuted-context $\alpha$-renaming (Probe~2) & $1000 / 1000$ & $100.00\%$ \\
Held-out lemma in retrieval shortlist (Probe~3) & $50 / 50$ & $100.00\%$ \\
\bottomrule
\end{tabular}
\end{center}

\noindent The compositional no-leakage certificate (Lemma~\ref{lem:comp-noleak})
holds empirically on real Lean proof states under the implemented parser.
Per \S\ref{sec:experiments}, this is necessary but not sufficient.

\subsubsection*{Head-to-head with variance bars}

We train all seven models (\S\ref{sec:baselines}) at seed~$0$. We then train
M3 and M5 at seeds~$1$ and~$2$ for variance bars on the principal comparison.

\begin{center}
\small
\begin{tabular}{lrrrrr}
\toprule
Model & params (M) & top-1 & top-5 & MRR (B) & AUC (C) \\
\midrule
M0b symbolic affordance-count MLP & $2.11$ & $0.359$ & $0.906$ & $0.001$ & $0.746$ \\
M1 byte-level text Transformer & $4.44$ & $0.382$ & $0.873$ & $0.001$ & $0.752$ \\
M2 weak hypergraph & $11.65$ & $0.364$ & $0.910$ & $0.002$ & $0.753$ \\
M3 strong role-aware hypergraph ($d = 448$) & $35.81$ & $0.382$ & $0.911$ & $0.002$ & $0.798$ \\
M4 role-typed GNN (grade-$0$/$1$ only) & $10.33$ & $0.362$ & $0.913$ & $0.002$ & $0.716$ \\
M5a structural-only RLIC ($\Kstruct$) & $36.04$ & $0.404$ & $0.911$ & $0.002$ & $0.822$ \\
M5 full RLIC ($\Kafford$) & $36.04$ & $0.401$ & $0.909$ & $0.002$ & $0.826$ \\
\bottomrule
\end{tabular}
\end{center}

\noindent Random baseline $= 1/9 \approx 0.111$; majority-class
(\texttt{simp}) baseline $\approx 0.25$. Both RLIC variants clear the
strongest non-RLIC baseline (M3) by approximately $2$~pp at single seed; M5
wins Task~C decisively. Averaged across three seeds each, the principal
comparison is:

\begin{center}
\small
\begin{tabular}{lcccc}
\toprule
Metric & M3 ($n{=}3$) & M5 ($n{=}3$) & $\Delta$ (M5 $-$ M3) & Significance \\
\midrule
top-1 & $0.3843 \pm 0.0031$ & $0.4037 \pm 0.0040$ & $+0.0194$ & $t \approx 6.7$, $p < 10^{-3}$ \\
top-5 & $0.9118 \pm 0.0007$ & $0.9103 \pm 0.0020$ & $-0.0015$ & tie (saturated) \\
AUC (C) & $0.8045 \pm 0.0084$ & $0.8309 \pm 0.0059$ & $+0.0264$ & $t \approx 4.4$, $p < 10^{-2}$ \\
MRR (B) & $0.0017 \pm 0.0001$ & $0.0015 \pm 0.0001$ & --- & both $\approx$ floor \\
\bottomrule
\end{tabular}
\end{center}

\noindent The principal comparison answers in the affirmative the empirical
question flagged in Remark~\ref{rem:hyperedge}: \emph{graded} cellular
structure with role-conditioned transport beats a parameter-matched strong
role-aware hyperedge baseline at LTM-Tiny scale on this subset.

\subsubsection*{Sliced evaluation}

Per-slice top-1 accuracy, averaged across three seeds:

\begin{center}
\small
\begin{tabular}{lrrr}
\toprule
Slice & M3 (mean $\pm$ s.d.) & M5 (mean $\pm$ s.d.) & $\Delta$ \\
\midrule
\texttt{rewrite\_heavy} & $0.334 \pm 0.020$ & $0.366 \pm 0.013$ & $+0.032$ \\
\texttt{multi\_arg\_app} & $0.386 \pm 0.008$ & $0.401 \pm 0.007$ & $+0.015$ \\
\texttt{induction} & $0.297 \pm 0.011$ & $0.315 \pm 0.015$ & $+0.018$ \\
\texttt{residual} & $0.407 \pm 0.002$ & $0.430 \pm 0.006$ & $+0.023$ \\
\bottomrule
\end{tabular}
\end{center}

\noindent The largest sliced top-1 gap is on \texttt{rewrite\_heavy}
($+3.19$~pp), consistent with the mechanism prediction of
\S\ref{sec:experiments}. M5 also wins the \texttt{residual} slice
($+2.35$~pp), so the gains are \emph{largest on} structurally complex
slices but \emph{not exclusive to} them; the strong reading of the targeted
hypothesis (``little or no benefit on structurally trivial states'') is not
supported in this run. The \texttt{constructor\_split} slice has fewer than
five states per seed in the filtered validation set and is suppressed.

\subsubsection*{Ablations on M5}

Each ablation in \S\ref{sec:baselines} applied to M5 at seed~$0$:

\begin{center}
\small
\begin{tabular}{lrrrr}
\toprule
Ablation & params (M, trainable) & top-1 & AUC (C) & $\Delta$ top-1 \\
\midrule
M5 (control) & $36.04$ & $0.4061$ & $0.8279$ & --- \\
(a) collapse all roles & $36.04$ & $0.3841$ & $0.8065$ & $-0.0220$ \\
(b) freeze $\rho^{r,\uparrow}, \rho^{r,\downarrow}$ to identity & $21.09$ & $0.3638$ & $0.7496$ & $-0.0423$ \\
(c) drop incidence-spectral PEs & $36.04$ & $0.4085$ & $0.8331$ & $+0.0024$ \\
(d) drop $\mathcal{N}_{\leftrightarrow}$ & $21.35$ & $0.3923$ & $0.8221$ & $-0.0138$ \\
\bottomrule
\end{tabular}
\end{center}

\noindent Ablation (c) is a control: incidence-spectral positional features
are not implemented in this prototype, so the perturbation is null up to
seed noise. The pattern decomposes the architectural contribution cleanly:

\begin{center}
\small
\begin{tabular}{lr}
\toprule
Increment & marginal top-1 gain \\
\midrule
Baseline (b): freeze transport, no roles & $0.364$ \\
$\quad+$ learnable role-conditioned transport, roles still collapsed (a) & $+2.03$~pp \\
$\quad+$ role labels (M5) & $+2.20$~pp \\
\midrule
Total architectural contribution & $+4.23$~pp \\
\bottomrule
\end{tabular}
\end{center}

\noindent The drop from ablation (a) is concentrated on slices where role
labels carry semantic content: $-2.2$~pp on \texttt{rewrite\_heavy}
($\role{source}$ vs $\role{target}$ on equalities), $-2.3$~pp on
\texttt{multi\_arg\_app} ($\role{argument}_k$ on applications), and exactly
$+0.1$~pp on \texttt{induction} (tie, within seed noise), which lacks
role-rich grade-$2$ structure. This is direct empirical evidence for the
role-labelling design choice. The drop from (b) is uniform across slices
($-2.6$ to $-4.9$~pp); without learnable role-conditioned transport the
model collapses to roughly the M2 weak-hypergraph level
($0.364 \approx 0.364$), validating Definition~\ref{def:transport} as the
mechanism that turns the four-complex hierarchy into a useful inductive
bias. Dropping $\mathcal{N}_{\leftrightarrow}$ costs $-1.4$~pp overall but
is approximately neutral on \texttt{rewrite\_heavy}; same-grade adjacency
is load-bearing for the \texttt{induction} and \texttt{multi\_arg\_app}
slices in particular.

\subsubsection*{Empirical caveats}\label{subsec:results-caveats}

We flag four points the reader should weigh against the headline.

\paragraph{Task~B is at floor for all models.}
With a $16{,}384$-premise pool, the retrieval head trained under InfoNCE
with $512$ random negatives reaches MRR $\approx 0.002$ and Recall@10
$\approx 0.003$ for every model, including the text Transformer and the
symbolic baseline. This is a head/loss issue (a small $\phi$ encoder
against a large premise vocabulary under random negatives is too weak a
learning signal), not an architectural one. The M5 vs M3 comparison on
Tasks~A and~C is not contaminated, but Task~B claims should be withdrawn
from this version and re-attempted with a larger $\psi$ encoder and harder
negative mining.

\paragraph{Task~C uses synthesised negatives.}
The LeanDojo benchmark records only successful tactic states; every
observed state is solvable by \emph{some} recorded tactic. We synthesise
$30\%$ negative states per split by swapping goals across theorems. This
yields the C\_AUC $\approx 0.83$ we report, but the negative-generation
procedure is a paper-specific choice that biases the value task toward
distinguishing structurally inconsistent context/goal pairs from coherent
ones. A negative pool mined from tactic-failure traces would be a stronger
evaluation; we leave this to follow-up.

\paragraph{Single subset, single split.}
We use the random split only. The \texttt{novel\_premises} split, which
holds out premises rather than theorems, is not exercised here.

\paragraph{No search loop.}
The $+1.94$~pp top-1 gain on Mode~1 tactic-family classification is
suggestive of gains in a full proof-search loop (MCTS or best-first), but
this is not demonstrated. MCTS coupling is the natural next empirical
question.

% =====================================================================
\section{System Overview}\label{sec:system}
% =====================================================================

\begin{small}
\begin{verbatim}
Lean proof state P
        |
        v
Parser (Algorithms 1-3 under Condition 3.6;
        Lemma 6.2 compositional certification;
        Algorithm 4 drift probes)
        |
        +--> K_struct(P)             pure structural state
        |
        +--> K_afford(P)             + state-only affordance cells
        |
        +--> K_a(P, A)               + candidate-action cells (Mode 2 only)
        |
        v
Role-conditioned incidence-transport layer
   (Definition 4.5; strict cellular-sheaf variant in the
    proof-state-cell-complex subcase, Definition 4.4)
        |
        +-------------------+-------------------+
        v                   v                   v
   policy head        value head         retrieval head
   (Mode 1: input     (input             (input
    K_afford)          K_afford)          K_afford)
        |                   |                   |
        +-------------------+-------------------+
                            |
                            v
                Lean kernel (trusted base)
\end{verbatim}
\end{small}

% =====================================================================
\section{Limitations}\label{sec:limitations}
% =====================================================================

\begin{itemize}
\item \textbf{Scope.} No MCTS, no recursive decomposition, no LM coupling, no Rocq backend, no proof-relevant equality.
\item \textbf{Empirical status.} Single LTM-Tiny run on a $45{,}450$-state filtered subset, three seeds on the principal M3-vs-M5 comparison, one seed elsewhere (\S\ref{subsec:results}). Findings: M5 beats the parameter-matched strong hyperedge baseline by $+1.94$~pp top-1 ($p < 10^{-3}$) and $+2.64$~pp Task-C AUC ($p < 10^{-2}$); largest sliced top-1 gap on \texttt{rewrite\_heavy} ($+3.19$~pp); Algorithm~\ref{sec:parser}.4 probes pass at $100\%$; ablations isolate learnable role-conditioned transport ($-4.23$~pp when frozen) and role labels ($-2.20$~pp when collapsed) as the two load-bearing components. Caveats: Task~B at floor for all models, Task~C uses synthesised negatives, single subset/split, no MCTS coupling. See \S\ref{subsec:results-caveats}.
\item \textbf{Parser configurability.} $\Pi$ is a hyperparameter.
\item \textbf{No-leakage probes are necessary, not sufficient.} Lemma~\ref{lem:comp-noleak} is the formal contract.
\item \textbf{Sheaf in the free case is a quiver representation.} Definition~\ref{def:sheaf} is coherence-constrained only after parser-imposed path identifications.
\item \textbf{The implemented model is not a sheaf.} Definition~\ref{def:transport} carries an independently parameterised downward map and no composition law.
\item \textbf{Lemma~\ref{lem:sens} vs typed hyperedges.} The higher-grade sensitivity result separates the model from grade-$1$-only baselines but not from typed-hyperedge baselines; \S\ref{sec:baselines} tests this empirically with a strong role-aware hyperedge baseline; \S\ref{subsec:results} resolves the empirical question in the affirmative at LTM-Tiny scale.
\item \textbf{Equivariance is conditional.} Theorem~\ref{thm:equiv} holds only under (i)--(vi) of \S\ref{sec:labels}; the retrieval caveat applies for retrieval encoders that depend on $K$.
\item \textbf{Mode 1 partly tests symbolic affordances.} Because affordance generators emit cells with names like \role{rewrite-opp}, \role{constructor-opp}, \role{induction-opp}, Mode~1 \emph{also} tests whether state-only symbolic affordances are predictive of tactic families, not solely whether the neural architecture learns proof structure from scratch. The symbolic affordance-count baseline in \S\ref{sec:baselines} is included to make this contribution structure visible. \S\ref{subsec:results} shows the symbolic baseline reaches top-1 $= 0.359$, while M5 reaches $0.401$ at single seed and $0.4037 \pm 0.004$ at three seeds: state-only affordance counts do account for a substantial fraction of the family-prediction signal, but not all of it.
\item \textbf{Mode 2 not exercised.} Requires a candidate-instance generator we do not provide.
\end{itemize}

% =====================================================================
\section{Future Work}\label{sec:future}
% =====================================================================

\begin{enumerate}
\item Parser-imposed path identifications enabling nontrivial cellular-sheaf coherence.
\item Identification of subclasses of RLICs that meet Definition~\ref{def:pscc}, supporting Hodge/sheaf Laplacians.
\item Family-universal candidate-instance generators for Mode~2.
\item MCTS coupling. The $+1.94$~pp Mode-1 top-1 gain of M5 over M3 at LTM-Tiny scale (\S\ref{subsec:results}) is suggestive of gains in a full proof-search loop; this is the most directly load-bearing follow-up empirical question.
\item Rocq instantiation; proof-relevant equality via path higher-grade cells.
\item Scaling study. We do not invoke a ``Large Topology Models'' label in this paper.
\item Stronger Task~B evaluation: larger $\psi$ encoder, hard negative mining against the corpus, and re-running over the \texttt{novel\_premises} split (the current run leaves Task~B at floor for all models; see \S\ref{subsec:results-caveats}).
\item Mined Task~C negatives from tactic-failure traces, replacing the cross-theorem goal-swap heuristic used here.
\end{enumerate}

% =====================================================================
\section{Conclusion}
% =====================================================================

We have given a role-labelled, graded incidence-complex representation of Lean proof states, a four-complex hierarchy separating structural state, state-only affordances, externally supplied candidate actions, and verified transitions, and an explicit no-leakage discipline with a computational boundary on the parser. Parser compliance is compositional: it reduces to per-generator compliance and deterministic truncation (Lemma~\ref{lem:comp-noleak}). The implemented model is a role-conditioned incidence-transport layer with explicit face-to-coface upward maps and coface-to-face downward maps, plus a class-pair-indexed same-grade map. The strict cellular-sheaf variant is identified precisely on the incidence category with face-to-coface arrows, with the free-incidence case acknowledged as a quiver representation. The architecture is equivariant under proof-complex isomorphism under stated conditions; degenerate cases reduce to set self-attention and to typed node--edge message passing; higher-grade cells admit parameter settings that distinguish complexes with identical grade-$1$ skeleta, though not those with identical typed-hyperedge representations. Evaluation protocol includes a weak hyperedge baseline, a strong role-aware hyperedge baseline, and a symbolic affordance-count baseline. The neural model never enters the trusted base; the Lean kernel does.

Narrowly stated: this specific role-labelled incidence-complex representation, the four-complex hierarchy of \S\ref{sec:hierarchy}, the no-leakage condition of \S\ref{sec:noleak} with its computational boundary and compositional certificate, the parser of \S\ref{sec:parser}, and the role-conditioned incidence-transport layer of \S\ref{sec:sheaves} and \S\ref{sec:arch}, together form a precise and defensible foundation for topology-aware neural theorem proving over Lean proof states.

\paragraph{Empirical outcome.} On a filtered LeanDojo Benchmark~4 subset at LTM-Tiny scale (\S\ref{subsec:results}; 45k train, 80k optimisation steps, three seeds on the principal comparison), the role-conditioned RLIC layer beats a parameter-matched strong role-aware hyperedge baseline by $+1.94$~pp top-1 ($t \approx 6.7$, $p < 10^{-3}$) and $+2.64$~pp Task-C AUC ($t \approx 4.4$, $p < 10^{-2}$); the largest sliced top-1 gap is on rewrite-heavy proof states ($+3.19$~pp), corroborating the mechanism prediction of \S\ref{sec:experiments}.7. Ablations isolate the learnable role-conditioned transport (Definition~\ref{def:transport}) and the role labels (Definition~\ref{def:roles}) as the two load-bearing architectural components, each contributing approximately half of the gain. All three Algorithm~\ref{sec:parser}.4 no-leakage probes pass at $100\%$ on $1{,}000$ real proof states. The principal limitation is the absence of a search loop: a gain on single-step Mode-1 prediction is suggestive of, but does not demonstrate, end-to-end theorem-proving improvement. MCTS coupling (\S\ref{sec:future}, item~4) is the natural next empirical question.

% =====================================================================
\begin{thebibliography}{99}
% =====================================================================

\bibitem{giusti2022simplicial}
L.\ Giusti, C.\ Battiloro, P.\ Di Lorenzo, S.\ Sardellitti, S.\ Barbarossa.
\emph{Simplicial Attention Neural Networks.}
arXiv:2203.07485.

\bibitem{ballester2024cellular}
R.\ Ballester, P.\ Hern\'andez-Garc\'ia, M.\ Bern\'aldez, S.\ Escalera, C.\ Casacuberta, D.\ Rojas.
\emph{Attending to Topological Spaces: The Cellular Transformer.}
arXiv:2405.14094.

\bibitem{bodnar2022sheaf}
C.\ Bodnar, F.\ Di Giovanni, B.\ P.\ Chamberlain, P.\ Li\`o, M.\ M.\ Bronstein.
\emph{Neural Sheaf Diffusion: A Topological Perspective on Heterophily and Oversmoothing in GNNs.}
arXiv:2202.04579.

\bibitem{hajij2022tdl}
M.\ Hajij, G.\ Zamzmi, T.\ Papamarkou, et al.
\emph{Topological Deep Learning: Going Beyond Graph Data.}
arXiv:2206.00606.

\bibitem{hansenghrist2019}
J.\ Hansen, R.\ Ghrist.
\emph{Toward a Spectral Theory of Cellular Sheaves.}
J.\ Appl.\ Comput.\ Topology 3, 2019. arXiv:1808.01513.

\bibitem{wang2017formulanet}
M.\ Wang, Y.\ Tang, J.\ Wang, J.\ Deng.
\emph{Premise Selection for Theorem Proving by Deep Graph Embedding (FormulaNet).}
arXiv:1709.09994.

\bibitem{paliwal2020graph}
A.\ Paliwal, S.\ Loos, M.\ Rabe, K.\ Bansal, C.\ Szegedy.
\emph{Graph Representations for Higher-Order Logic and Theorem Proving.}
AAAI 2020.

\bibitem{blaauwbroek2024graph2tac}
L.\ Blaauwbroek, M.\ Ol\v{s}\'ak, J.\ Rute, F.\ Massolo, J.\ Piepenbrock, V.\ Pestun.
\emph{Graph2Tac: Online Representation Learning of Formal Math Concepts.}
arXiv:2401.02949.

\bibitem{yang2023leandojo}
K.\ Yang, A.\ M.\ Swope, A.\ Gu, R.\ Chalamala, P.\ Song, S.\ Yu, S.\ Godil, R.\ Prenger, A.\ Anandkumar.
\emph{LeanDojo: Theorem Proving with Retrieval-Augmented Language Models.}
arXiv:2306.15626.

\bibitem{xin2024deepseek15}
H.\ Xin, Z.\ Z.\ Ren, J.\ Song, et al.
\emph{DeepSeek-Prover-V1.5.}
arXiv:2408.08152.

\bibitem{deepseekv2}
DeepSeek-AI.
\emph{DeepSeek-Prover-V2.}
arXiv preprint (verify bibliographic details against the current arXiv listing before submission).

\bibitem{rocqdocs}
The Rocq Prover documentation.
\emph{Typing rules --- Calculus of Inductive Constructions.}
\url{https://rocq-prover.org}.

\end{thebibliography}

\end{document}
